% Source: https://artofproblemsolving.com/community/c3h1347165_area_of_the_kite
% Stable identity: aops-topic-1347165-q01
The diagonals of a kite are perpendicular. The diagonal of length $8$ cuts the kite into two triangles with this diagonal as common base; their heights are the two parts of the other diagonal, which add up to $6$.

\begin{asy}
usepackage("amsmath");
size(4.5cm);
defaultpen(fontsize(10pt));
path ra(pair P,pair u,pair v,real s=0.2){return P+s*unit(u)--P+s*unit(u)+s*unit(v)--P+s*unit(v);}
pair P=(0,6), Q=(3,0), R=(0,-2), T=(-3,0), O=(0,0);
draw(P--Q--R--T--cycle,linewidth(0.9));
draw(P--R,blue+dashed); draw(T--Q,red+dashed);
draw(ra(O,Q-O,P-O,0.35));
draw(P--(-4.3,6),gray+dotted); draw(R--(-4.3,-2),gray+dotted);
draw(T--(-3,-3.3),gray+dotted); draw(Q--(3,-3.3),gray+dotted);
draw((-4,-2)--(-4,6),blue,Arrows(TeXHead)); draw((-3,-3)--(3,-3),red,Arrows(TeXHead));
label("$8$",(-4,2),W,blue); label("$6$",(0,-3),S,red);
label("\mbox{}",(min(currentpicture,user=true)+max(currentpicture,user=true))/2,black); // keeps the text after the figure black when LaTeX typesets the labels inline
\end{asy}

So the area is
\[
\dfrac{1}{2}\cdot8\cdot h_1+\dfrac{1}{2}\cdot8\cdot h_2=\dfrac{1}{2}\cdot8\cdot6=\boxed{24},
\]
and since $24$ lies in the allowed range $24$ to $36$, \textbf{yes}, Lou can fly the kite.
